

\begin{frame}{Compositional data}
\only<1>{
\note{
IFirst of all we should know what simplex space is, I will frequently refer to simplex space. 
A simplex space is  collection of vectors with non-negative values that always sum to a constant.
This space is known as the \emph{simplex} and it is common to set \( k = 1 \).
We refer to data from this space as \emph{compositional data}.

The problem I'm going to talk about is how can we apply a random forest to this space and why is this important?
}


\vspace{1em}

A simplex is a convex space of all vectors with non-negative components that sum to a constant. 



\begin{center}
\colorbox{blue!10}{\parbox{0.8\textwidth}{
\begin{center}
\[
S^c_k = \left\{ (y_1, \dots, y_k)^\top \in \mathbb{R}^k \,\middle|\, y_i \geq 0, \ \sum_{i=1}^k y_i = \textcolor{blue}{c} \right\}
\]
\end{center}
}}
\end{center}



 It is common to set \(  \textcolor{blue}{c = 1} \). We refer to data from this space as \emph{compositional data}.
 

}



\only<2>{

\textbf{A practical view of compositional data}

\vspace{0.3cm}

Let $\mathbf{Y} = (Y_1, Y_2, \dots, Y_K)$ be a vector with positive components.

\vspace{0.3cm}

The compositional transformation is defined as:
\[
\mathbf{S} =
\left(
\frac{Y_1}{\sum_{i=1}^K Y_i}, \dots,
\frac{Y_K}{\sum_{i=1}^K Y_i}
\right)
\]

\vspace{0.3cm}

This vector lies in the simplex space.

\vspace{0.5cm}

However, in this transformation we lose information about the \textbf{absolute scale}, and we cannot uniquely recover the original vector.

But we preserve the \textbf{relative information}.
In particular, for any two components $i$ and $j$:
\[
\frac{Y_i}{Y_j} = \frac{S_i}{S_j}
\]


\note{
Improved spoken version

We can also define compositional data in a more practical way.

Suppose we have a vector with non-negative components and no other restrictions.

If we divide each component by the total sum, we obtain a new vector that belongs to the simplex space.

This new vector is a compositional representation of the original data.

An important point is that we cannot recover the original vector uniquely from its compositional form.

In this transformation, we lose information about the absolute scale of the data.

However, what we can fully preserve is the relative information.

In particular, the ratios between components remain exactly the same after normalization.

So for any two components, the ratio is identical in both representations.

It is commonly assumed in this field that absolute values are not important.

The only meaningful information in compositional data is the relative information between components.
}


}

% \only<3>{
% \vspace{.2cm}
%    \textbf{Example}: consider 
% \(
% \mathbf{Y} = (2000, 3000, 5000),
% \)
% which represents the amount of money (in dollars) allocated to \textit{stocks}, \textit{bonds}, and \textit{cash} in an investment portfolio.  

% Then the following vector is a compositional data:
% \begin{align*}
% \scalebox{1.2}{$\left(\frac{2000}{2000+3000+5000},\;\frac{3000}{2000+3000+5000},\;\frac{5000}{2000+3000+5000}\right)$} = (0.2, 0.3, 0.5)
% \end{align*}

% \medskip

% Multiply by 100:
% \[
% 100 \cdot \mathbf{s}(\mathbf{Y}) = (20, 30, 50)
% \]

% \medskip

% Same proportions, different scale (percentages).
% \note{
% Here we see a portfolio allocation example.

% After normalization, we obtain proportions of total investment.

% Multiplying by 100 expresses them as percentages, which are more intuitive in finance.

% The relative allocation (20\%, 30\%, 50\%) remains unchanged regardless of the total invested amount.
% }
% }


\only<3>{
\note{Confirm:In this slide, consider a job application pool of 50 applicants. Among them, 10 have a mathematics background, 15 have a statistics background, and 25 have a computer science background.

If we normalize this vector by dividing each component by the total sum, we obtain a compositional vector.

Each component corresponds to the proportion of applicants with a different academic background.

By multiplying this vector by 100, we obtain percentages. For example, 20% of the applicants have a mathematics background.

The important point is that the information remains the same, but the scale changes—from proportions to percentages.

In general, multiplying a compositional vector by any positive constant does not change the relative information; it only changes the scale.

Now, this simple normalization step looks simple, but it has a very important statistical consequence.
}

Consider a job application pool of $50$ applicants classified by academic background.
\vspace{0.3cm}

\begin{columns}
    \begin{column}{0.4\textwidth}
        \[
        \mathbf{Y} = (10, 15, 25)
        \]
    \end{column}
    \begin{column}{0.5\textwidth}
        \begin{tabular}{lr}
            Mathematics & 10 \\
            Statistics  & 15 \\
            Computer Science & 25 \\
        \end{tabular}
    \end{column}
\end{columns}

\vspace{0.3cm}
We normalize this vector to obtain a compositional representation:
\vspace{0.2cm}

\begin{columns}
    \begin{column}{0.5\textwidth}
        \[
        \mathbf{s}(\mathbf{Y}) = \left(\frac{10}{50}, \frac{15}{50}, \frac{25}{50}\right)
        \]
    \end{column}
    \begin{column}{0.4\textwidth}
        \[
        = (0.2, 0.3, 0.5)
        \]
        Interpreted as proportions.
    \end{column}
\end{columns}

\vspace{0.2cm}
Multiplying by $100$ gives percentages:
\vspace{0.2cm}

\begin{columns}
    \begin{column}{0.5\textwidth}
        \[
        100 \cdot \mathbf{s}(\mathbf{Y}) = (20, 30, 50)
        \]
    \end{column}
    \begin{column}{0.4\textwidth}
        \textbf{Key idea:} The information is identical — only the scale changes.
    \end{column}
\end{columns}


}

\end{frame}


% \begin{frame}{Statistical problems with compositional data}
% 1)spurouis corelation
% 2)The (sample) covariance matrix of the data is singular with rank equal to $D - 1$, since
% \[
% \sum_{i=1}^{D} \operatorname{cov}(y_1, y_i)
% = \operatorname{cov}\left(y_1, \sum_{i=1}^{D} y_i \right)
% = \operatorname{cov}(y_1, 1)
% = 0,
% \]
% or equivalently,
% \[
% \sum_{i=2}^{D} \operatorname{cov}(y_1, y_i)
% = -\operatorname{var}(y_1),
% \]
% where each $y_i$ is a component of the composition $\mathbf{y}$.
% 3) Evidently, the problems and issues associated
% with compositional data are too numerous to list here and this is beyond the scope of
% this review paper.

    
% \end{frame}

% \begin{frame}{What methodology is more approperiate?}
% Aitchison suugest to tranfer data...

% dicson had another idea...

% a long argue staeted between them ... and finally Aitchison used stupid word and dicson lift the discussion for evere...
% Aitchison keep continue work on his idea and bring a lot of attention this his idea so many researcher for a while tried to his idea ... 
% aclsy and .. belive this made bad effect on the the field ... and for a while many researcher tried to say you have to use transformation methods...

% They belived 
% \end{frame}









